\chapter{绪论}
\section{研究背景与问题}
椭圆型偏微分方程常用于描述稳态扩散和势场。对于由不同材料组成的区域，扩散系数在材料界面发生跳跃，真解通常仅在每个子区域内光滑。即使解连续，梯度也可能不连续；在给定非零跳跃数据时，解本身亦可不连续。因此，单纯在整个区域内使用一个光滑近似并忽略界面条件，不能完整表达所要求解的问题。

传统差分或有限元方法通过网格及界面处理建立离散系统。随机神经网络方法则预先生成隐藏层参数，把近似写为固定特征的线性组合，只求输出系数。对于线性PDE，内部方程和线性边界、界面条件均可转化为线性约束。这为CPU上的可重复小规模研究提供了条件，但隐藏特征的相关性、配置点覆盖及约束尺度会影响求解结果。

本文研究的问题是：在固定随机特征和明确计算预算下，如何评价内部与界面配置点的选取及分配？残差较大的位置增加点数是否必然改善解误差？当系数对比度提高时，误差、矩阵条件性和计算成本是否具有一致的变化趋势？这些问题不能仅用训练损失回答，需要独立真解验证与合理静态基线。

\section{相关研究及贡献边界}
Sun、Dong和Wang将局部随机网络与间断Galerkin框架结合，提出多个格式并给出带假设的误差分析\cite{dgpreprint}。其中弱形式积分、面耦合及近似条件与本文强形式配置不同，不能将两者统称为同一离散格式。Li和Wang使用不同子域的随机网络求解界面问题，采用中心差分近似特征导数，并在随机点上叠置内部、界面和边界约束\cite{interfacepreprint}。其预印本总结明确把自适应配置点、初始化范围及数值分析列为后续方向；这为本文研究动机提供依据，却不保证本文方案在现阶段具有新颖性。

Dang、Wang和Jiang的AG-RaNN使用频率信息、残差和增长策略构造新增特征，并涉及域分裂\cite{ag}。本文阅读的是2025年的第三版预印本。其特征或网络架构适应与本文固定特征下的空间点适应不同；本文不直接套用其在特定正则性、参数及采样假设下的理论结论。Shang等结合单位分解、PCA与重叠Schwarz预条件来改善局部随机特征系统求解\cite{schwarzpreprint}，提示条件性不仅与点选取有关，也与特征冗余和求解器有关。

残差驱动采样本身已有明确先行研究。Wu等比较固定采样及RAD、RAR-D等方法\cite{rar}。较近的Residual-Christoffel方法依据算子作用后的特征空间几何构造密度、权重和白化\cite{christoffel}，其中“残差空间”并不等于当前近似解残差的幅值。本文的简单联合指标与这些技术的差异尚不足以支持原创算法声明。RNN-PG 2023论文正式书目信息已核实\cite{pg}，但全文尚未获得，因此本文不依据其未读推导建立算法或归纳作者结论。

\section{研究内容及组织}
本文的工作包括四个部分。第一，独立实现Poisson五点差分和随机特征配置，验证数学数据与导数。第二，构造零与非零跳跃的双子域制造解，研究加权最小二乘、尺度平衡与特征参数。第三，在共享特征和明确预算下比较四种采样及消融，保留无效、恶化和局部峰精度不足的结果。第四，给出标准有限维稳定性结论的自含推导，并通过扰动及花形界面近似复现讨论数值限制。

第二章介绍模型、制造解、差分和最小二乘；第三章推导局部随机特征离散系统；第四章说明采样及预算机制；第五章报告真实数值实验；第六章总结有限范围结论与未解决问题。附录给出补充模型与复现路径。

\section{研究范围与学术真实性}
本文不要求原创定理或重大算法创新。实验均为程序实际输出，原文报告数值与本项目数据分别列示。所有比较至少五种子，代表图不依据最好误差挑选。正式版和预印本存在题名、年份差异时分别记录；无法核实的全文内容不补写。本文尚未开展真实外部数据、三维、多层增长或完整Schwarz实现，也不作优于PINN或传统界面有限元的直接数值结论。
